\section{Método} \label{sec:method}
    O método utilizado para analisar os artigos consiste na 
    análise do \textit{abstract} e da conclusão de cada 
    artigo de interesse, com o objetivo de identificar 
    pontos relevantes como a tecnologia proposta, o tipo
    de jogo testado, o aspecto que se buscou aprimorar 
    no jogo e o objetivo proposto no artigo.


\subsection{Base de dados usada na busca}
    Considerando que os jogos constituem uma área 
    multidisciplinar \cite{mortazavi2024dynamic}, 
    foi utilizado o mecanismo de busca Google Scholar,
    um mecanismo de busca on-line gratuito que permite 
    pesquisar literatura acadêmica e científica, 
    incluindo artigos revisados por pares, teses, 
    livros e decisões judiciais.


\subsection{Termos de busca}
    Os termos de busca utilizados foram definidos 
    com base na proposta inicial de explorar DDA 
    em jogos, sendo eles:

    \begin{enumerate}
        \item \textit{Intelligent Adaptive Games}
        \item \textit{Fighting Games DDA}
    \end{enumerate}

    A busca na literatura incluiu artigos completos de 
    periódicos e conferências publicados em inglês
    no período entre 2020 e 2026.


\subsection{Critérios de inclusão e exclusão}
    Os artigos selecionados tiveram que satisfazer os 
    seguintes critérios para serem incluídos na análise:

    \begin{itemize}
        \item Data de publicação entre 2020 e 2026 
        \item 5 ou mais citações
    \end{itemize}

    Com base nesses critérios, foram incluídas publicações
    em periódicos e conferências que desenvolveram e 
    avaliaram empiricamente técnicas para o ajuste 
    dinâmico de dificuldade ou o ajuste adaptativo de 
    níveis de jogo em entretenimento digital.


\subsection{Processo de triagem}
    Com a aplicação dos termos de busca 
    estabelecidos, a seleção foi feita por meio da 
    análise do \textit{abstract} e da conclusão de cada
    artigo, a fim de determinar sua relevância para a 
    proposta. 
    
    Após essa avaliação inicial, restaram seis 
    publicações que satisfizeram as condições de 
    inclusão para uma análise mais detalhada e coleta de
    dados.


